% !TeX root = ../main.tex
\section{Conclusions}
\label{sec:conclusions}

We derived a closed-form finite-deformation Biot coefficient for a saturated
solid with logarithmic skeleton and mineral laws and plastic distention.
A scalar equation determines mineral volume, and the coefficient follows
directly from that volume, total volume, and pore pressure.  The expression
recovers \(B_0=1-K/K_s\) at the reference state and connects the pressure
contribution to total stress with the stress that drives plastic flow.

Plastic distention carries loading history into the coefficient.  The
material-point calculations show a different unloading response from that
of the virgin material at the same total deformation.  At nonzero pressure,
the change in coefficient feeds back into the driving stress and restrains
dilation for the chosen energy and flow law.  The implicit return mapping
updates these quantities together, and automatic differentiation retains
their coupled changes in the Newton Jacobian.  Independent checks of mass,
energy-derived stress, flow, yield consistency, and derivatives verify the
update on the smooth Drucker--Prager branch.

The coupled finite-element calculations verify the elastic specialization
against the Mandel pressure and horizontal-displacement solution.  A larger
compression test shows how the coefficient varies across the specimen during
drainage.  Mesh and time-step comparisons measure numerical sensitivity;
these calculations provide implementation verification for the chosen laws.
The next steps are to verify coupled flow with plastic deformation,
including its contribution to water storage, and to calibrate and test the
model against material data.
